\section{Discrete and series-parallel boundaries}
\label{sec:a-boundaries}

Independent finite resistance choices behave differently from intervals.
For pressure, two cycles in one block already permit a Subset-Sum
encoding. For an individual arc, classical electrical current signs
preserve interval-hull extrema on series-parallel graphs, so the first
hard block rank is three. At unbounded block rank, fixed nominations
still admit high-precision arc optimization, but exact arithmetic and
nomination uncertainty impose separate limits. Finally, a worst-case
hull theorem does not decide whether one finite resistance choice
realizes a prescribed flow.

Throughout this section, series-parallel means that the undirected graph
has no $K_4$ minor; equivalently, its biconnected components admit
series-parallel decompositions
\citep{cosmellopez2017-k4}. No terminals are designated globally. Except in the
explicit continuous-law statements, all laws are $g_e(x)=\beta_e x|x|$,
all input data are rational, and every allowed resistance is positive.
Finite sets are listed explicitly. Scenario domains contain all passive
states for the allowed nominations and resistances. Capacities are
subsequent tests, not constraints filtering that domain.

\subsection{A pressure maximum encoding Subset Sum}
\label{subsec:a-bound-pressure}

\begin{theorem}[Discrete pressure hardness at block rank two]
\label{thm:a-bound-pressure}
Maximizing a prescribed potential difference over independent two-point
resistance sets is $\NP$-hard on simple graphs of maximum degree three
with one rank-two block and one bridge. Nominations are fixed small
integers. The reduction has a rational separation gap of polynomial
encoding length. It excludes an algorithm polynomial in input length
and accuracy bits unless $\mathrm P=\NP$. With unnormalized positive
integer resistances, even absolute-error-one value approximation is
$\NP$-hard. These are not strong-hardness claims.
\end{theorem}
\begin{proof}
Take positive Subset-Sum integers $a_1,\ldots,a_n,K$, put $S=\sum_i a_i$,
and preprocess $K>S$ directly. Use vertices $0,1,2,3$ and arcs
\[
 (0,2),\ (2,1),\ (0,3),\ (3,1)
 \quad\text{of resistances}\quad
 \tfrac12,\ \tfrac16,\ \tfrac16,\ \tfrac12.
\]
Replace $(2,3)$ by a directed path of $n$ edges, whose internal
nominations vanish, with choices
\[
 \beta_i\in\left\{\frac1{2n},\frac1{2n}+\frac{a_i}{2K}\right\}.
 \qquad
 \tau=\sum_i\beta_i=\frac12+\frac{\sum_i a_i\sigma_i}{2K},
 \quad \sigma_i\in\{0,1\}.
\]
The resulting theta block has rank two. First set the nominations at
$0,1,2,3$ to $(4,-4,3,-3)$. The unique positive root $q$ of
\begin{equation}
 (12\tau+1)q^2+10q-23=0
 \label{eq:a-bound-theta}
\end{equation}
lies in $(0,2)$: the polynomial is negative at zero, positive at two,
and strictly increasing on $[0,\infty)$. Assign cross-path flow $q$ and,
in the displayed order, outer flows
\[
 \frac{q+1}{2},\quad\frac{7-q}{2},\quad
 \frac{7-q}{2},\quad\frac{q+1}{2}.
\]
Their divergences are the stated nominations. The two outer path drops
agree by symmetry; their intermediate potential difference is
$[23-10q-q^2]/12=\tau q^2$. Thus every cycle equation holds, and
\cref{thm:a-pre-existence} identifies this as the physical state.
The common outer drop is
\[
 \pi_0-\pi_1=\frac{(q+1)^2}{8}+\frac{(7-q)^2}{24}
             =2+\frac{(q-1)^2}{6}.
\]
Add a leaf $4$ and an arc $(4,1)$ of resistance three. Give $4$
nomination one and change $b_1$ to $-5$. Its flow is one, and block
aggregation preserves the theta state. Hence
\begin{equation}
 F:=\pi_4-\pi_0=1-\frac{(q-1)^2}{6}\le1,
 \label{eq:a-bound-peak}
\end{equation}
with equality exactly when $q=1$, or equivalently $\tau=1$ and
$\sum_i a_i\sigma_i=K$. The fixed nominations are $(4,-5,3,-3,1)$.
The graph is simple, has maximum degree three, and has just the asserted
blocks. Fixed yes/no instances of Subset Sum, for example $(a_1,K)=(1,1)$
and $(2,1)$, handle preprocessing within the same graph class.

If a choice does not sum to $K$, then $|\tau-1|\ge1/(2K)$.
Subtracting the polynomial at $1$ from its value at $q$ gives
\[
 |q-1|=\frac{12|\tau-1|}{(12\tau+1)(q+1)+10}
       \ge\frac6{31K+18S}.
\]
Here $q<2$ and $\tau\le(K+S)/(2K)$ bound the denominator by
$31+18S/K$. Consequently every non-target choice satisfies
\begin{equation}
 F\le1-\Delta,\qquad \Delta=\frac6{(31K+18S)^2}.
 \label{eq:a-bound-discrete-gap}
\end{equation}
An absolute-error estimate below $\Delta/3$ distinguishes the two
optimum values. Its requested accuracy has polynomial binary length.
A near-optimal discrete choice with error below $\Delta$ must itself
encode a target subset in a yes instance, which can be checked by
integer arithmetic.

Multiply every resistance by $6nK$. The four fixed coefficients become
$3nK,nK,nK,3nK$, the path choices become
$\{3K,3K+3na_i\}$, and the bridge coefficient becomes $18nK$.
Flows do not change and potentials scale by $6nK$. Multiplying again by
$T=(31K+18S)^2$ gives an integer peak $6nKT$ and a gap
$6nKT\Delta=36nK\ge36$. An absolute-error-one estimate compared with
the peak minus $18$ decides Subset Sum. These coefficients have
polynomial binary length, but need not have polynomial numerical
magnitude; this scaling proves neither strong nor relative-error
hardness.
\end{proof}

The binary path sum is the usual Subset-Sum mechanism. The network
feature is the strict interior pressure maximum in
\cref{eq:a-bound-peak}. In particular, the individual outer and cross
flows in this gadget are monotone in $\tau$; pressure hardness alone
does not establish arc hardness. The interval version of this same
rank-two problem remains covered by \cref{thm:a-blk-joint}.
Existing discrete sizing hardness has a different scope:
\citet[Theorem 4.7]{robinius2019-discrete-sizing} treats minimum-cost
arc sizing subject to potential feasibility, already on trees, and
credits the earlier water-network hardness of Yates, Templeman, and
Boffey. Here resistance choices are adversarial parameters for one
unfiltered pressure objective, without a cost budget or potential
bounds. On a tree with fixed nominations this objective is linear in
the resistances. Thus that design result neither proves nor conflicts
with the rank-two boundary above.

\subsection{Finite comparison, electrical signs, and resistance hulls}
\label{subsec:a-bound-hulls}

\begin{lemma}[Finite secant comparison]
\label{lem:a-bound-secant}
On any connected graph, let $x,\pi$ and $y,\rho$ be the passive states
for the same nomination and two families of continuous strictly
increasing laws $g_e,G_e$ vanishing at zero. Set
\[
 r_e=\frac{G_e(y_e)-G_e(x_e)}{y_e-x_e}\quad(y_e\ne x_e),
 \qquad d_e=G_e(x_e)-g_e(x_e),
\]
and choose any $r_e>0$ when $y_e=x_e$. For a target arc $a=(u,v)$,
let $j$ be the unit electrical flow from $u$ to $v$ with resistances $r$.
Then
\begin{align}
 (\rho_u-\rho_v)-(\pi_u-\pi_v)&=j^\top d,\label{eq:a-bound-secant-pressure}\\
 y_a-x_a&=\frac{\sum_{e\ne a}j_ed_e-(1-j_a)d_a}{r_a}.
 \label{eq:a-bound-secant-flow}
\end{align}
Moreover, $0\le j_a\le1$.
\end{lemma}
\begin{proof}
Strict increase makes every secant positive, including the arbitrary
choice multiplying a zero flow difference. Subtracting the physical
systems gives
\[
 A(y-x)=0,\qquad A^\top(\rho-\pi)=\operatorname{diag}(r)(y-x)+d.
\]
Write $j=\operatorname{diag}(1/r)A^\top h$ with
$Aj=e_u-e_v$. Pairing the second identity with $j$ makes its first
right-hand term $h^\top A(y-x)=0$, proving
\cref{eq:a-bound-secant-pressure}. Taking its $a$ coordinate then gives
\cref{eq:a-bound-secant-flow}. The electrical maximum principle puts
$h_u\ge h_v$, so $j_a\ge0$; decomposition of the acyclic unit electrical
flow into source-to-sink paths gives $j_a\le1$, as in the proof of
\cref{lem:a-cac-lipschitz}.
\end{proof}

\begin{lemma}[Classical adjacent-terminal signs]
\label{lem:a-bound-signs}
In a series-parallel graph, the sign of every unit electrical current
between the endpoints of an existing target edge is independent of all
positive electrical resistances. Edges off the target block have zero
current. The signs are computable by one rational electrical solve
with unit resistances.
\end{lemma}
\begin{proof}
Block decomposition makes every off-block current zero. The target
block is a bridge or is two-terminal series-parallel with the target
endpoints as terminals, by
\citet[Lemma 9]{eppstein1992-parallel-recognition-of-series-parallel}.
In a two-terminal decomposition, series components transmit the same
positive current and parallel components split it into positive
currents. Induction fixes each edge's direction. The bridge case is
immediate. A reduced rational Laplacian solve at unit resistances has
polynomial bit complexity and recovers these signs, including zeros.
\end{proof}

This is classical confluence theory: \citet[Theorems 0 and 1]{duffin1965-topology-of-series-parallel-networks}
relates resistance-independent current directions to graph structure,
and his Section 5 treats monotone nonlinear resistor characteristics.
The restriction to adjacent terminals is essential. Arbitrary terminal
pairs in a theta block need not have these fixed signs.

\begin{theorem}[Continuous-law envelopes]
\label{thm:a-bound-envelope-structural}
Fix a balanced nomination on a series-parallel graph. For each edge
let $\Theta_e$ be a nonempty compact parameter space, with a jointly
continuous law $g_e(t,\theta)$ such that each member is strictly
increasing in $t$ and vanishes at zero. Edge choices are independent.
For each target arc, one deterministic envelope network attains its
maximum flow over all scenarios; a second attains its minimum.
The two envelope profiles are independent of the nomination.
\end{theorem}
\begin{proof}
Define $g_e^{\max}(t)=\max_{\theta\in\Theta_e}g_e(t,\theta)$ and
$g_e^{\min}(t)=\min_{\theta\in\Theta_e}g_e(t,\theta)$.
Compactness and joint continuity give continuous envelopes and attained
extrema. For $s<t$, evaluate the maximum at a member attaining the
maximum at $s$, and the minimum at a member attaining the minimum at
$t$; strict increase proves strict increase of both envelopes. They
vanish at zero, so \cref{thm:a-pre-existence} applies even if they are
bounded.

Let $\sigma_e$ be the signs in \cref{lem:a-bound-signs}. For maximization
choose $G_a=g_a^{\min}$ on the target, choose $G_e=g_e^{\max}$ when
$e\ne a$ and $\sigma_e>0$, and choose $G_e=g_e^{\min}$ when
$e\ne a$ and $\sigma_e<0$. On zero-current edges choose any member.
Compare its physical flow $y$ with any original flow $x$ using
\cref{eq:a-bound-secant-flow}. Every numerator term is nonnegative,
so $y_a\ge x_a$. At $y_e$ choose a member attaining the selected
pointwise envelope value, independently on every edge. The complete
state $y,\rho$ satisfies those original laws, hence is their unique
state. The upper bound is therefore attained. Interchange minima and
maxima for the lower bound. Neither envelope construction used $b$.
\end{proof}

The theorem allows coupling among the parameters of one edge; only
independence between edges is required. It is structural for general
continuous families: without an input representation it implies no
algorithm. The fixed envelope flow is continuous in $b$: bounded nominations give
bounded acyclic flows, and summing the continuous edge drops along a
spanning tree bounds normalized potentials. Every subsequential limit
satisfies the limiting state equations and hence is unique. Therefore,
for any nonempty compact balanced nomination set $\mathcal U$ it gives
\[
 \max_{b\in\mathcal U,\,\theta\in\prod_e\Theta_e}x_a(b,\theta)
 =\max_{b\in\mathcal U}x_a^{\rm env}(b).
\]
A realizing member profile may depend on $b$. Different target arcs
usually require different envelope networks. Neither a simultaneous
vector optimum nor preservation of existential feasibility follows.
The pressure identity \cref{eq:a-bound-secant-pressure} also gives a
cactus version: for any terminal pair, orient the unit electrical flow
along its block path. Bridge currents and the positive split between
the two paths in every traversed cycle have fixed signs. Choosing the
upper envelope on positive-current edges and the lower on negative-current
edges, with the same attainment argument, maximizes that potential
difference. Reverse the choices to minimize it.

\begin{corollary}[Separate monotonicity and exact finite-set validation]
\label{cor:a-bound-finite-positive}
On a series-parallel graph, a target arc flow is separately monotone
in each scalar parameter that changes only one law affinely, provided
all intermediate complete laws are admissible. Thus products of compact
scalar sets with attained endpoints and admissible interval hulls have
the same target-flow extrema as their interval hulls. This equality
continues to hold jointly over any compact nomination set.

For quadratic laws, finite rational resistance sets, a balanced rational
nomination box, and fixed maximum block rank, exact arc extrema and
robust capacity validation are polynomial-time on series-parallel graphs.
In particular this holds whenever $\blockrank\le2$.
\end{corollary}
\begin{proof}
Write the changing law as $g_e(t,\theta)=g_e^0(t)+\theta f_e(t)$.
If $f_e(x_e(\theta_0))=0$ at some parameter, the entire physical state
at $\theta_0$ satisfies every parameter value; uniqueness makes it
constant. Otherwise $f_e(x_e(\theta))$ has a fixed strict sign by
continuity. State continuity follows directly from bounded flows,
locally uniform convergence of the laws, and uniqueness: along a
convergent parameter sequence, any flow subsequence has a convergent
subsequence, normalized potentials are bounded by summing edge drops along
a spanning tree, and its limit satisfies the limiting equations.
The flow bound is the fixed total positive nomination, by acyclicity.

For $\theta'<\theta$ compare the two states by the secant lemma. Only
$d_e=(\theta-\theta')f_e(x_e(\theta'))$ is nonzero. Its coefficient is
$j_e/r_a$ for $e\ne a$ and $-(1-j_a)/r_a$ for $e=a$. Its weak sign
is fixed by \cref{lem:a-bound-signs,lem:a-bound-secant}. This proves
monotonicity without smoothness or a nonzero derivative assumption.
Move an optimizer one coordinate at a time to an appropriate endpoint;
its objective stays optimal and the final parameters belong to the
original compact sets. The same argument holds with an optimizing
nomination fixed. Compactness and the continuity just proved give the
needed optimizer.

For the computational claim, replace each resistance set by its endpoint
interval and use \cref{thm:a-blk-arc}. An exact finite resistance optimizer
can also be recovered. Compute the optimum $V$, fix one resistance to
its lower endpoint, and reoptimize the remaining interval problem. Keep
that endpoint if its exact optimum equals $V$, and otherwise fix the
upper endpoint. Separate monotonicity ensures one endpoint retains $V$.
Polynomially many fixed-rank optimizations and algebraic comparisons
suffice. At the end the fixed-resistance algorithm returns an algebraic
nomination optimizer; no rational exact physical state is promised.

Finally, a $K_4$ minor must occur within a block. One way to see this is
to take its subdivision: as $K_4$ has degree three, prune each connected
minor branch set to a tree joining its at most three attachment points;
the three paths meet at one vertex, giving a subdivision. This
biconnected subdivision lies in one block. Cycle rank cannot increase
under edge deletion or contraction, while $K_4$ has rank three.
Thus $\blockrank\le2$ excludes such a minor.
\end{proof}

\subsection{An additive algorithm without a block-rank bound}
\label{subsec:a-bound-envelope-algorithm}

\begin{theorem}[Fixed-nomination series-parallel arc optimization]
\label{thm:a-bound-envelope-algorithm}
Fix rational balanced nominations on a series-parallel graph, and allow
independent resistance intervals or explicitly listed finite positive
resistance sets. For any target arc and rational $\eps>0$, its maximum
and minimum have rational enclosing intervals of width at most $\eps$,
and allowed rational endpoint scenarios attaining each extremum within
$\eps$ can be computed in time polynomial in input length and
$\log^+(1/\eps)$. Block rank is unrestricted.
\end{theorem}
\begin{proof}
Write $L_e,U_e$ for the resistance endpoints. The pointwise maximum of
$L_et|t|$ and $U_et|t|$ uses $U_e$ for $t\ge0$ and $L_e$ for $t\le0$;
the minimum reverses them. Thus \cref{thm:a-bound-envelope-structural}
constructs one asymmetric quadratic network for each extremum, with
positive rational coefficients $c_e^+,c_e^-\in\{L_e,U_e\}$. Its state
$x^*$ minimizes
\begin{equation}
 E(x)=\frac13\sum_e\bigl[c_e^+(x_e^+)^3+c_e^-(x_e^-)^3\bigr],
 \qquad Ax=b,
 \label{eq:a-bound-cubic-energy}
\end{equation}
where $x_e^+=\max\{x_e,0\}$ and $x_e^-=\max\{-x_e,0\}$.

This is a rational second-order cone program of linear size. Introduce
nonnegative $p_e,n_e$ with $x_e=p_e-n_e$; at a minimum both cannot be
positive, since subtracting their minimum preserves flow and decreases
energy. For $z\ge0$, the cubic epigraph $z^3\le u$ has the lift
\[
 z^2\le w,\qquad w^2\le uz,\qquad u,w,z\ge0.
\]
The first two constraints are rotated second-order cones. They imply
$z^4\le uz$, and conversely $w=z^2$ is feasible whenever $u\ge z^3$;
the zero case is immediate. Apply this to every $p_e,n_e$.
We next give explicit bit bounds, since a conic representation alone
does not guarantee the required output accuracy.

If $b=0$, all flows vanish. Otherwise put
$B=\sum_v|b_v|>0$, $\beta_L=\min_e L_e$, and $\beta_U=\max_e U_e$.
The passive state has $|x_e^*|\le B/2$. Choose a rational tree-supported
flow $x^0$ with $|x_e^0|\le B/2$, and a fundamental-cycle matrix
$C\in\{0,1,-1\}^{E\times k}$ whose chord submatrix is the identity.
Every feasible flow is $x^0+Cz$ and $z^*$ is the vector of chord flows,
so $|z_i^*|\le B/2$. When $k=0$, return the tree flow. Otherwise
minimize $f(z)=E(x^0+Cz)$ on $[-B,B]^k$. On this cube,
\[
 \begin{aligned}
 |x_e^0+(Cz)_e|&\le R:=(m+1)B,\\
 |\partial_i f|&\le G_0:=m\beta_U R^2,
 &\|\nabla f\|_2&\le G:=1+mG_0.
 \end{aligned}
\]
Rational evaluation and gradients use the sign of each affine flow to
select its cubic piece, and have polynomial bit cost. A coordinate
outside the cube yields a separating coordinate hyperplane. At a point
with $f(z)>t$, its rational gradient separates that point from
$\{f\le t\}$. A zero gradient there certifies that sublevel empty.

For any $\delta>0$, a sublevel with $t\ge f(z^*)+\delta$ contains the
ball about $z^*$ of radius
\[
 r_\delta=\min\{B/4,\delta/(2G)\}.
\]
The cube lies in the ball of radius $mB$. Thus the rational ellipsoid
weak-feasibility algorithm, using the above exact separation oracle,
finds a feasible point whenever the sublevel has this ball, in a number
of bit operations polynomial in the input and
$\log((1+mB)/r_\delta)$. Threshold bisection between zero and $f(0)$,
with inner accuracy $\delta/4$, returns a rational $\widehat z$ with
$f(\widehat z)\le f(z^*)+\delta$. More explicitly, a failed ball-search
at threshold $t$ implies $t<f(z^*)+\delta/4$, while a successful search
supplies a point with value at most $t$. These lower and upper bounds
allow bisection to terminate with objective error $\delta$.
All ellipsoid arithmetic uses polynomial-length rounding, as in the
standard rational weak-optimization theorem
\citep[Theorem 2.5.9]{dadush2012-integer-programming}.
Its bounded-body hypotheses are met by the cube; if a globally
Lipschitz objective is used in that formulation, extend each scalar
primitive beyond $[-R,R]$ by its tangent affine functions. The resulting
convex rational oracle agrees on the entire cube and has the same
gradient bound. The returned flow $\widehat x=x^0+C\widehat z$
satisfies conservation exactly.

It remains to convert energy error to flow and scenario error. For an
envelope law $G_e$, the function $G_e(t)-\beta_Lt|t|$ is nondecreasing.
The scalar inequality
\[
 (u|u|-v|v|)(u-v)\ge\tfrac12|u-v|^3
\]
follows by factoring on each same-sign half-line and by
$u^2+v^2\ge(|u|+|v|)^2/2$ for opposite signs. Integrating this
monotonicity along a segment, and using
$G(x^*)^\top(\widehat x-x^*)=0$, gives
\begin{equation}
 E(\widehat x)-E(x^*)\ge
 \frac{\beta_L}{6}\sum_e|\widehat x_e-x_e^*|^3.
 \label{eq:a-bound-energy-error}
\end{equation}
Hence energy error $\delta=\beta_L\eta^3/6$ gives
$\|\widehat x-x^*\|_\infty\le\eta$, and
$[\widehat x_a-\eta,\widehat x_a+\eta]$ encloses the exact extremum.

Choose each original resistance endpoint using the sign of
$\widehat x_e$ and the envelope rule; either endpoint is allowed at
zero. A wrong sign choice can occur only when $|x_e^*|\le\eta$.
The selected law $g_e'$ consequently has residual
\[
 |g_e'(x_e^*)-G_e(x_e^*)|\le\beta_U\eta^2.
\]
Let $x'$ be its true passive flow. Pair the two potential systems with
the circulation $x'-x^*$ and apply the same scalar inequality to $g'$.
With $D=\|x'-x^*\|_\infty$ this gives
\[
 \frac{\beta_L}{2}D^3\le m\beta_U\eta^2D,
 \qquad D\le\sqrt{\frac{2m\beta_U}{\beta_L}}\,\eta.
\]
Set $C_0=1+2m\beta_U/\beta_L$ and
$\eta=\min\{\eps/2,\eps/C_0\}$. The interval then has width at most
$\eps$ and the selected allowed scenario has target flow within
$\eps$ of optimal. All internal tolerances have polynomial binary
length. No exact sign decision for an irrational flow, or numerical
solver oracle, is assumed.
\end{proof}

The convex energy principle and weak convex optimization are classical.
The quantitative conclusion combines them with an envelope construction
and control of endpoint selection near zero flow. We do not assign
priority to the envelope identity itself: nonlinear circuit tolerance
analysis predates this formulation, and the directly relevant
Hasler--Wang 1993 paper \emph{Parameter tolerances in non-linear resistive
circuits: worst case analysis based on monotonicity}
\citep{hasler1993-parameter-tolerances} has not been read in full.
We draw no theorem-level priority conclusion from it. The claims here
are the explicitly proved model and output guarantees.

More specifically, \citet[Theorem 3]{chauffoureaux1990-monotonicity}
gives a structural source-to-resistor monotonicity criterion through a
linear circuit test, motivates endpoint evaluation for worst-case
analysis, and sketches extensions to resistor parameters and the
observed branch itself. Our finite secant argument accommodates zero
derivatives and merely continuous laws, while the quadratic algorithm
specifies binary input, accuracy bits, and an allowed endpoint scenario.
\citet[Sections 4--5 and 8--9 of the report version]{zemanian1999-ratings}
composes finite rated characteristic segments on series-parallel
networks, tests existence of a within-ratings operating point, and
recovers a designated branch by choosing the reduction order. That is
a deterministic rated-characteristic problem; the uncertainty here
instead varies independently selected complete constitutive laws.
The report's formulas, tables, and interpolation do not provide the
binary-precision guarantee proved above.

\subsection{Arc hardness at block rank three}
\label{subsec:a-bound-arc-hardness}

A high-resistance probe converts a pressure gap into a rational arc
threshold. We first record the monotonicity that controls its feedback
on the original network. If a probe $(s,t)$ carries flow $z$, its
removal changes the old nomination to $b-z e_s+z e_t$. Let $F(z)$ be
the old-network pressure $\pi_s-\pi_t$ at this nomination. For two
distinct values $z,w$, strict increase of every law gives
\begin{align*}
 &\bigl(b(z)-b(w)\bigr)^\top\bigl(\pi(z)-\pi(w)\bigr)\\
 &\hspace{1cm}=\sum_e\bigl(g_e(x_e(z))-g_e(x_e(w))\bigr)
                         \bigl(x_e(z)-x_e(w)\bigr)>0.
\end{align*}
The left side is $-(z-w)(F(z)-F(w))$, so $F$ is strictly decreasing.
A probe of resistance $M$ therefore obeys $F(z)=Mz|z|$; its sign is
the sign of $F(0)$, and it reduces the absolute terminal pressure.
All these comparisons concern the unique full physical state, whose
existence was already proved.

\begin{theorem}[Discrete arc classification]
\label{thm:a-bound-arc-hardness}
At fixed maximum block rank, deciding whether a finite resistance choice
and a nomination in a bounded rational balanced box can give $x_a>c$
is in $\NP$; deciding whether every scenario satisfies $x_a\le c$ is
in $\coNP$. Both are complete for these classes already with fixed
small integer nominations, two-point resistance choices, and a simple
biconnected graph of maximum degree three and rank three. The same
completeness holds for existence of $x_a\ge c$ and its complementary
universal strict inequality. All thresholds are rational. The reduction
has a polynomially encoded positive gap, excluding polynomial-time
optimization in accuracy bits unless $\mathrm P=\NP$.
\end{theorem}
\begin{proof}
Start with \cref{thm:a-bound-pressure}, before integer scaling, and set
\[
 D=\left\lceil\frac{10000}{\Delta}\right\rceil,
 \qquad H=1-\Delta/2,\qquad M=HD^2,\qquad c=1/D.
\]
Add the target edge $(4,0)$ of resistance $M$. The former leaf and this
edge form an ear between $1$ and $0$ in the theta block. The entire
graph is biconnected, has rank three, and still has maximum degree three.
For each discrete choice the old pressure is
$F(0)=1-(q-1)^2/6\in(5/6,1]$.
The new edge flow $z$ is positive, and monotonicity gives
\[
 0<Mz^2=F(z)\le1,\qquad z\le2/D<1,
\]
since $H\ge1/2$. The old nomination changes by $z e_0-z e_4$.
Both nominations lie in a box with total absolute coordinate bound at
most $18$. The fixed path $4,1,2,0$ has resistance sum $11/3$.
Thus \cref{lem:a-cac-lipschitz} gives, uniformly in all uncertain
cross-path resistances,
\begin{equation}
 0\le F(0)-F(z)\le264z\le528/D<\Delta/8.
 \label{eq:a-bound-probe-loss}
\end{equation}
For a target subset, $F(0)=1$, whence $Mz^2>1-\Delta/8>H=Mc^2$
and $z>c$. For every non-target choice,
$Mz^2\le1-\Delta<H$ and $z<c$. These strict gaps prove the claimed
reductions for both weak and strict existential comparisons, without
changing equality conventions. Their complements give the two universal
versions. If all edges require capacities, bounds $[-16,16]$ on every
other edge are harmless by the total-nomination flow bound; the target
can have lower bound $-16$.

For a target choice the pressure excess over $H$ is greater than
$3\Delta/8$; for a non-target choice the deficit is at least
$\Delta/2$. Since $z+c\le3/D<4/D$ and $M\le D^2$,
\begin{equation}
 |z-c|=\frac{|F(z)-H|}{M(z+c)}
       \ge\frac{3\Delta}{32D}.
 \label{eq:a-bound-arc-gap}
\end{equation}
An additive estimate of error below $\Delta/(64D)$ separates yes and
no instances. Near-optimal discrete scenarios also identify target
subsets at this accuracy. Every number has polynomial binary length.
Writing $A_0=31K+18S$, multiplication of all resistances by
$6nKA_0^2$ makes them positive integers, including the probe coefficient
$6nK(A_0^2-3)D^2$. It changes no flow or capacity. In particular,
resistance scaling does not amplify \cref{eq:a-bound-arc-gap}.

For membership, guess one input resistance index per edge. With these
coefficients fixed, \cref{thm:a-blk-arc} computes the exact target
extrema over the continuous balanced nomination box in polynomial bit
time at fixed block rank. Compare the algebraic maximum with $c$ using
the required weak or strict inequality. A verifier therefore needs no
rational nomination or physical-state witness. For several capacities,
guess an offending edge and direction as well. Complementation gives
universal membership. Together with the rank-three reduction this proves
completeness. It does not give a pressure-threshold membership theorem
for sums over arbitrarily many blocks.
\end{proof}

\Cref{cor:a-bound-finite-positive,thm:a-bound-arc-hardness} yield a sharp
block-rank boundary: finite-resistance robust arc validation is polynomial
at rank at most two and $\coNP$-complete when rank three is permitted.
This does not say every rank-three graph is hard; series-parallel graphs
can have arbitrarily large block rank and still satisfy
\cref{thm:a-bound-envelope-algorithm} for fixed nominations.
The reduction is from Subset Sum with a small flow gap. It establishes
neither strong hardness nor fixed-error hardness under the stated
fixed-nomination normalization.

\subsection{The universal graph class for individual-arc hulls}
\label{subsec:a-bound-characterization}

\begin{theorem}[Universal arc-hull characterization]
\label{thm:a-bound-characterization}
For a finite connected simple graph $G$, the following are equivalent:
\begin{enumerate}
\item $G$ is series-parallel.
\item For every fixed balanced nomination and positive quadratic
resistance interval family, each target arc flow is separately monotone
in every resistance coordinate.
\item For every fixed balanced nomination and independent compact
positive resistance sets with attained endpoints, each target-flow
maximum and minimum equals its interval-hull counterpart.
\end{enumerate}
If $G$ has a $K_4$ minor, failure can be witnessed by rational data,
small integral nominations, and one two-point uncertain resistance.
Given a $K_4$ subdivision, the witness has polynomial encoding length
and an explicit rational gap.
\end{theorem}
\begin{proof}
The positive implications are \cref{cor:a-bound-finite-positive}.
For the converse, take the theta in \cref{eq:a-bound-theta}.
Keep nominations $(4,-4,3,-3)$ and a
single cross edge of resistance $\theta$. Its reverse pressure is
\[
 F_\theta(0)=\pi_1-\pi_0=-2-(q-1)^2/6.
\]
At $\theta=1$ it equals $-2$. At
$\theta_L=23/108$ and $\theta_U=71/12$, substitution in
\cref{eq:a-bound-theta} gives $q=3/2$ and $q=1/2$, respectively,
and both pressures equal $-49/24$.

Add the target edge $(1,0)$ of resistance $M=10^8$. This is exactly
$K_4$. Its flow $z<0$ obeys $F_\theta(z)=Mz|z|$ with old nominations
$b+z e_0-z e_1$. The preceding probe monotonicity gives
\[
 F_\theta(0)<F_\theta(z)<0,\qquad
 Mz^2<-F_\theta(0)\le49/24<4.
\]
Thus $|z|<2/\sqrt M<1$. Use the nomination bound $B=18$ and path
$1,2,0$ of resistance sum $2/3$ in \cref{lem:a-cac-lipschitz}:
\[
 |F_\theta(z)-F_\theta(0)|\le48|z|
 <96/\sqrt M<1/96.
\]
The interior probe pressure is greater than $-2$, whereas both endpoint
pressures are less than $-49/24+1/96=-65/32$. Thus the interior target
flow is strictly larger than both endpoint flows. More quantitatively,
its pressure advantage exceeds $1/32$, and all three probe flows have
magnitude below $2/\sqrt M$. Since they are negative,
\[
 F_I-F_E=M(z_I-z_E)(|z_I|+|z_E|),\qquad
 z_I-z_E>g:=\frac1{128\sqrt M}=\frac1{1280000}.
\]
Here $I$ denotes the interior scenario and $E$ either endpoint.

Every $K_4$ minor contains a subdivision, by the branch-tree argument
in \cref{cor:a-bound-finite-positive}. Embed the above six-edge model
on that subdivision $H$. Distribute each fixed branch resistance
equally over its path, and set every internal nomination to zero.
On a subdivided uncertain path use a fixed total
$d=\theta_L/2$ on all but one edge, distributed equally, and resistance
$\theta-d$ on the remaining edge. If the path has length one, use
$d=0$. All endpoints are positive. Series aggregation preserves the
three states and their target gap; choose any actual edge of the probe
path as target.

It remains to restore all edges of $G$ outside $H$ without silently
deleting them. Put zero nominations at extra vertices and give every
extra edge the same positive integer resistance $R$. Let $m=|E(G)|$.
For each of the three scenarios, the theta flow with $q=1$, zero probe
flow, and zero extra-edge flows is conservation-feasible. Its energy is
$(10+\theta)/3<6$, independently of subdivisions. The complete physical
minimum consequently has every extra-edge flow bounded by
$u=(18/R)^{1/3}$. Its total positive nomination is seven, so every
physical edge flow has magnitude at most seven.

Restriction of that state to $H$ is the passive state at induced
balanced nominations $b'=A_Hx_H$. They obey
$|b'_v|\le7\deg_H(v)$ and
$\|b'-b\|_1\le2mu$. The original nominations obey the same coordinate
bounds, so a containing box has total absolute bound at most $14m$.
For the actual target edge of resistance $\beta_a$, apply
\cref{lem:a-cac-lipschitz} to the single-edge terminal path, and use
$\bigl|s|s|-t|t|\bigr|\ge |s-t|^2/2$. This gives
\[
 |x'_a-x_a|^2\le\frac2{\beta_a}|\Delta(\pi_u-\pi_v)|
 \le56m\|b'-b\|_1\le112m^2u.
\]
Take
\[
 R=18\left(\frac{2000m^2}{g^2}\right)^3.
\]
It is an integer, has polynomial binary length, and gives
$|x'_a-x_a|<g/4$ in each scenario because $112/2000<1/16$.
The restored interior flow still exceeds both restored endpoint flows
by more than $g/2$. Every other resistance remains fixed, so this
violates both separate monotonicity and the two-point hull equality.
All path splits also have polynomial rational encoding length.
The existential characterization needs only some sufficiently large
$R$; the displayed bound proves the stronger encoding claim once a
subdivision has been selected.
\end{proof}

\subsection{Exact arithmetic and nomination uncertainty}
\label{subsec:a-bound-unbounded}

The additive guarantee in \cref{thm:a-bound-envelope-algorithm} cannot
be silently upgraded to exact threshold comparison. This distinction
already occurs without uncertainty.

\begin{theorem}[A probe-cactus arithmetic equivalence]
\label{thm:a-bound-arc-srs}
Exact weak lower comparison of one quadratic arc flow with a rational
threshold is $\SRS$-hard on simple series-parallel graphs of maximum
degree three, with fixed unit source--sink nominations and fixed positive
integer resistances. The threshold $1/2$ suffices for hardness.
On the restricted family consisting of a two-terminal cactus branch
with two terminal bridges in parallel with one probe edge, weak lower
comparison of the designated probe flow oriented from the unit source
to the unit sink is polynomial-time many-one equivalent to $\SRS$.
\end{theorem}
\begin{proof}
The integer cactus encoder in \cref{thm:a-cac-srs} produces a chain of triangles that is simple and has maximum degree
three, with unit-demand effective
resistance $D>0$ and positive integer threshold $H$, such that
\[
 \sum_i\sqrt{a_i}\le K\quad\Longleftrightarrow\quad D\ge H.
\]
Its terminal degrees are two. Add new terminals $S,T$, each joined to
an old terminal by a unit-resistance bridge, and add the probe $(S,T)$
of resistance $H+2$. The cactus branch has effective resistance $D+2$.
The graph is simple, maximum degree three, and two-terminal
series-parallel. Unit demand splits into positive branch flows $1-x$
and $x$; homogeneity gives
\[
 (D+2)(1-x)^2=(H+2)x^2,
 \qquad x=\frac{\sqrt{D+2}}{\sqrt{D+2}+\sqrt{H+2}}.
\]
Thus $x\ge1/2$ if and only if $D\ge H$, including equality.
No square root is evaluated when constructing the instance. For
preprocessed trivial cases use a triangle whose direct resistance is
two and whose two-edge path has coefficients one and one. Its effective
resistance is $1/2$; after two unit terminal bridges the branch has
resistance $5/2$. A probe of resistance one yields a fixed yes instance,
and resistance nine a fixed no instance, in the same literal family.

Conversely, let $R>0$ be the effective resistance of any such cactus
branch and $\beta>0$ the probe resistance. The probe flow is strictly
between zero and one. Thresholds $c\le0$ are trivial yes instances and
$c\ge1$ are trivial no instances. For $0<c<1$,
\[
 x\ge c\quad\Longleftrightarrow\quad
 R\ge\beta\left(\frac c{1-c}\right)^2.
\]
This is precisely the rational weak lower cactus-pressure comparison
reduced to $\SRS$ in \cref{thm:a-cac-srs}. The equivalence does not
interchange weak inequalities with strict capacity violations.
\end{proof}

The construction merges the growing number of cactus cycles into one
block after adding the probe, so it does not contradict the exact
fixed-block-rank algorithm. As stated in \cref{subsec:a-pre-computation},
$\SRS$ is not known to have a polynomial-time algorithm or to be
$\NP$-hard. \Cref{thm:a-bound-arc-srs} is an arithmetic classification
of the specified family, not an ordinary $\NP$-hardness result or an
upper bound for all series-parallel graphs.

Allowing a nomination box creates a different obstruction. The pressure
hardness used next is due to Th\"urauf; only its quantitative transfer
to an arc objective is needed here.

\begin{theorem}[Nomination hardness on series-parallel graphs]
\label{thm:a-bound-nomination}
With fixed positive rational quadratic resistances, maximizing one
signed arc flow over a balanced rational nomination box is $\NP$-hard
on series-parallel graphs. Robust upper-capacity validation is
$\coNP$-hard. The reduction has a positive rational gap of polynomial
encoding length, so a polynomial algorithm in input length and accuracy
bits would imply $\mathrm P=\NP$. Block rank grows with the instance;
no $\NP$ or $\coNP$ membership claim is made for this class.
\end{theorem}
\begin{proof}
We use the pressure instance of
\citet[Figure 2, equation (3), and Lemmas 4.3 and 4.17]{thurauf2022-deciding-the-feasibility-of-a}.
For positive Partition numbers $S_1,\ldots,S_n$ with $n\ge3$ and
$K=\sum_iS_i$, it has arcs $(s,z_i^+),(z_i^+,t),(z_i^+,z_i^-)$,
with respective resistances $1/S_i^2,1/S_i^2,1$. Its signed nomination
box, intersected with balance, is
\[
 b_s\in[0,K/2],\quad b_t\in[-K/2,0],\quad
 b_{z_i^+}\in[0,S_i],\quad b_{z_i^-}\in[-S_i,0].
\]
Its total absolute coordinate bound is $3K$. The main block is
$K_{2,n}$, and all $(z_i^+,z_i^-)$ edges are pendant bridges. Define
\begin{align*}
 \eps_0&=1-\left(1-\frac1{8K^2}\right)^2,\\
 M_0&=\max\{1-\eps_0+\eps_0^2,\ 1-\eps_0^2/K^2\},\qquad
 \eps_1=(1-M_0)/5,\\
 T&=\max\{1-\eps_1^2/(K^2n^2),\ M_0+4\eps_1\},
 \qquad \gamma=1-T>0.
\end{align*}
These are rational numbers of polynomial encoding length, with
$0<T<1$. Th\"urauf proves that a yes Partition instance admits a
nomination with $\pi_s-\pi_t\ge1$, while in a no instance every
nomination has $\pi_s-\pi_t<T$. The cited statements concern his
pressure subproblem (2), which imposes conservation and passive laws
but no potential bounds. Thus no feasibility filtering is imported.
Small Partition instances can be decided directly and replaced by a
fixed yes or no instance with at least three positive items.

Add a direct probe $(s,t)$ with
\[
 D=\left\lceil\frac{1000K(3K+2)}{\gamma}\right\rceil,
 \qquad H=1-\gamma/2,\qquad M=HD^2,\qquad c=1/D.
\]
It is one more parallel branch, so the graph remains series-parallel;
the main block rank increases from $n-1$ to $n$. For an original
nomination $b$, let $F_b(z)$ be the old pressure at
$b-z e_s+z e_t$. The probe relation is $F_b(z)=Mz|z|$ and $F_b$ is
strictly decreasing. If $F_b(0)\le0$, then $z\le0<c$. Otherwise
$F_b(z)\le F_b(0)$. A two-edge $s$--$t$ path has resistance sum at most
two, and the uniform flow bound $3K$ gives
$|F_b(0)|\le18K^2$. For a yes witness,
\[
 0<z\le6K/D<1.
\]
The original and induced nominations lie in an enlarged balanced box
with absolute coordinate bound at most $3K+2$. By
\cref{lem:a-cac-lipschitz} on that two-edge path,
\[
 0\le F_b(0)-F_b(z)\le8(3K+2)z
 \le\frac{48K(3K+2)}D<\gamma/4.
\]
The yes witness therefore has $F_b(z)>1-\gamma/4>H$ and $z>c$.
In a no instance, each positive probe flow satisfies
$F_b(z)\le F_b(0)<T<H$, so $z<c$; nonpositive flows satisfy this too.
The auxiliary perturbed nominations need not retain the entry/exit
sign restrictions; the sensitivity estimate is valid on their enlarged
signed box, and the source theorem is applied only to original
nominations.

The yes-case pressure margin exceeds $\gamma/4$. Rationalizing the
probe equation and using $z+c\le(6K+1)/D$ gives the flow gap
\[
 g_{\rm arc}=\frac{\gamma}{4(6K+1)D}.
\]
For a positive no-case flow, $z<2/D$ and $H-F_b(z)>\gamma/2$, so
$c-z>\gamma/(6D)\ge g_{\rm arc}$. A nonpositive flow has deficit
at least $1/D$. Hence the optimum is above $c+g_{\rm arc}$ in a yes
instance and below $c-g_{\rm arc}$ in a no instance. All encoding
lengths, including that of the required approximation error, are
polynomial. Other capacities may be set to $[-3K,3K]$ and do not
restrict the original scenarios.
\end{proof}

This is a consequence of Th\"urauf's Partition pressure construction,
not a new pressure-hardness source. It is compatible with
\cref{thm:a-bound-envelope-structural}: eliminating resistance uncertainty
leaves a nomination optimization that can still be hard. The gap does
not establish strong hardness or rule out running time polynomial in
reciprocal error. Related robust-design characterizations must also be
read with their computational scope intact: Th\"urauf, Gr\"ubel, and
Schmidt \citep{thurauf2025-adjustable-robust-nonlinear-network-design}
reduce fixed-design robust-feasibility checking to polynomially many
worst-case nonlinear subproblems and use them in an adversarial design
algorithm. A polynomial number of such subproblems is not a polynomial
time bound for solving them, or for the full design problem.

\subsection{Prescribed-flow realization and existential capacity design}
\label{subsec:a-bound-realization}

\begin{theorem}[An inverse-design quantifier boundary]
\label{thm:a-bound-realization}
For independent quadratic resistance intervals on any graph, realizing
a prescribed rational flow is rational linear feasibility and is
polynomial-time. For finite resistance sets, prescribed rational flow
realization is $\NP$-complete, even on one simple cycle with fixed
nominations $(2,-2)$, unit prescribed flow in an acyclic orientation,
positive integer resistances, and two choices per uncertain edge.
Existence of a finite resistance choice satisfying unit arc capacities
is also $\NP$-complete on the fixed-nomination single-cycle class.
\end{theorem}
\begin{proof}
For a target $\bar x$, first test $A\bar x=b$. Put
$t_e=\bar x_e|\bar x_e|\in\Q$. The remaining interval problem is
\[
 A^\top\pi=\operatorname{diag}(t)\beta,\qquad L\le\beta\le U.
\]
After fixing a potential reference, these are rational linear
constraints. Zero target flows impose zero endpoint potential drop and
cause no difficulty. A feasible rational solution of polynomial length
can be recovered by linear programming. With finite sets, guessing one
resistance index per edge leaves rational edge drops, whose membership
in $\operatorname{im}A^\top$ is checked by rational linear algebra.
This proves $\NP$ membership for prescribed-flow realization on any
graph, without a physical-state algebraicity assumption.

For hardness take positive Subset-Sum integers $a_1,\ldots,a_n,K$,
$n\ge2$, $S=\sum_i a_i$, and $0<K\le S$. Join source to sink by a
long directed path of $n$ edges and one direct edge. Give the long edges
choices $\{1,1+a_i\}$ and the direct edge fixed resistance $D=n+K$.
Prescribe flow one on every arc and nominations $2,-2$ at the terminals,
zero elsewhere. This is a simple cycle, maximum degree two, with an
acyclic arc orientation. The long-path total resistance is
\[
 \theta=n+\sum_i a_i\sigma_i.
\]
The two path drops at the target flow are $\theta$ and $D$, so the
conserved target is physical exactly when the chosen subset sums to
$K$. Trivial source cases use a triangle: both long-edge option sets
$\{1,2\}$ and direct resistance three give a yes instance; both sets
$\{1,3\}$ and direct resistance three give a no instance.

For any resistance choice, every long edge has the same physical flow
$p$, and the direct edge has flow $z$, with $p+z=2$. Both are strictly
positive because their path drops agree and their total flow is
positive. Unit upper capacities imply $p,z\le1$, hence force $p=z=1$.
Conversely the target satisfies every capacity. Absolute capacities
give the same equivalence. For general fixed-nomination single-cycle
inputs, guess the resistance choices and verify all capacities by the
exact scalar piecewise-quadratic cycle algorithm, or by
\cref{thm:a-blk-arc}. Thus this existential capacity problem is in $\NP$
on the stated class. No general-rank capacity-design membership claim
is needed.
\end{proof}

The two-path quadratic split gives an explicit approximation gap:

\[
 p=\frac{2\sqrt D}{\sqrt\theta+\sqrt D},\qquad z=2-p,
 \qquad |p-1|=\frac{|D-\theta|}{(\sqrt D+\sqrt\theta)^2}.
\]
If no subset sums to $K$, then $|D-\theta|\ge1$. With
$M_1=n+S+K$, both coefficients are at most $M_1$, so
$|p-1|\ge1/(4M_1)$. The minimum maximum arc load is therefore one
in a yes instance and at least $1+1/(4M_1)$ in a no instance.
The minimum maximum deviation from the all-one target is zero or at
least $1/(4M_1)$, respectively. Both objectives are convex functions of
the flow, but minimizing them over this finite scenario set is hard in
accuracy bits. This is a precision-dependent gap, not fixed-error or
strong hardness under the fixed-small-nomination normalization.

Replacing $\{1,1+a_i\}$ by $[1,1+a_i]$ lets $\theta$ range throughout
$[n,n+S]$. Every nontrivial instance then realizes the target because
$D=n+K$ lies in this interval. Moreover, the all-low and all-high finite
scenarios have $p\ge1$ and $p\le1$, respectively; since the complete
flow vector is affine in $p$, their segment contains the all-one target.
It can thus belong both to the interval flow region and to the convex
hull of finite-scenario flows while belonging to no finite scenario.
Containment of every scenario in a convex capacity set is preserved by
convexification; intersection with that set is not. This explains why
robust validation on a cycle is polynomial by
\cref{cor:a-bound-finite-positive}, while existence of a capacity-feasible
finite design is $\NP$-complete.
